% =============================================================================
%  Section 2.1 — Các đặc tính của nghệ thuật Chèo
% =============================================================================

\section{Các đặc tính của nghệ thuật Chèo}

Chèo là một loại hình sân khấu dân gian truyền thống của Việt Nam, mang tính tự sự cao và phản ánh đời sống xã hội thông qua hình thức diễn xướng tổng hợp. Tuy nhiên, trong phạm vi nghiên cứu này, điều quan trọng không chỉ nằm ở giá trị văn hóa của chèo, mà ở các đặc trưng cấu trúc và biểu đạt của nó – những yếu tố trực tiếp chi phối cách thức mô hình hóa tri thức dưới dạng bản thể luận (Ontology) và Knowledge Graph.

Trước hết, đặc trưng nổi bật nhất của chèo là cấu trúc tự sự phân lớp. Một vở chèo không được tổ chức theo chuỗi hồi – cảnh với quan hệ nhân quả chặt chẽ như kịch phương Tây, mà bao gồm nhiều lớp diễn tương đối độc lập về nội dung. Mỗi lớp diễn thường tập trung vào một tình huống, một xung đột hoặc một khía cạnh tính cách nhân vật cụ thể. Điều này tạo ra một cấu trúc dạng mô-đun, trong đó các lớp diễn có thể được xem như các đơn vị tri thức riêng biệt, vừa có tính độc lập tương đối, vừa có thể liên kết với nhau thông qua cốt truyện tổng thể. Chính đặc điểm này đặc biệt phù hợp với việc biểu diễn tri thức dưới dạng đồ thị, nơi mỗi lớp diễn có thể được mô hình hóa như một thực thể (entity) hoặc một nút (node), và các mối quan hệ giữa chúng được biểu diễn dưới dạng các cạnh (edges).

Thứ hai, hệ thống nhân vật trong chèo mang tính tượng hình và có cấu trúc phân loại rõ ràng. Các mô hình nhân vật như Đào, Kép, Hề, Lão và Mụ không chỉ đóng vai trò là các nguyên mẫu (archetype) trong nghệ thuật biểu diễn, mà còn tạo ra một hệ phân cấp (hierarchy) tự nhiên, rất thuận lợi cho việc xây dựng các lớp (class) trong Ontology. Từ các mô hình cơ bản này, có thể phát triển các tiểu loại chi tiết hơn dựa trên đặc điểm tính cách, vai trò trong cốt truyện hoặc vị thế xã hội. Điều đó cho thấy hệ thống nhân vật chèo không chỉ giàu tính biểu đạt mà còn có cấu trúc phân loại rõ ràng, phù hợp với các phương pháp biểu diễn tri thức dựa trên phân cấp và kế thừa.

Thứ ba, mối quan hệ giữa các thành phần trong chèo mang tính ngữ nghĩa phong phú và đa dạng. Các quan hệ như “nhân vật tham gia lớp diễn”, “lớp diễn thuộc vở diễn”, hay “nhân vật đối lập/ hỗ trợ lẫn nhau” đều có thể được xác định một cách rõ ràng và nhất quán. Đây là cơ sở quan trọng để xây dựng các quan hệ (relations) trong Knowledge Graph, cho phép thực hiện các truy vấn phức tạp và suy luận ngữ nghĩa. Đặc biệt, do mỗi lớp diễn thường làm nổi bật một khía cạnh cụ thể của nhân vật, nên mối quan hệ giữa nhân vật và lớp diễn có thể được khai thác để biểu diễn sự phát triển hoặc biến đổi của nhân vật trong toàn bộ vở chèo.

Thứ tư, nội dung của chèo thường dựa trên các mô-típ quen thuộc của văn học dân gian, với hệ thống giá trị đạo đức được thể hiện tương đối rõ ràng. Điều này giúp việc trích xuất và chuẩn hóa tri thức trở nên thuận lợi hơn, do các khái niệm và quan hệ thường mang tính ổn định và có thể được định nghĩa một cách nhất quán. Đồng thời, sự xuất hiện của các nhân vật có tính chất trung gian hoặc phức tạp về tâm lý trong một số vở chèo cũng mở ra khả năng biểu diễn các quan hệ ngữ nghĩa tinh vi hơn, chẳng hạn như các mức độ mâu thuẫn nội tâm hoặc chuyển biến tâm lý.

Từ các đặc trưng trên, có thể thấy rằng nghệ thuật chèo sở hữu một cấu trúc tri thức tự nhiên phù hợp với việc mô hình hóa dưới dạng Ontology và Knowledge Graph. Do đó, trong khóa luận này, hệ thống Ontology được thiết kế theo hướng phân tầng, bao gồm các lớp chính như vở diễn, lớp diễn và nhân vật, cùng với các quan hệ ngữ nghĩa tương ứng. Cách tiếp cận này không chỉ phản ánh trung thực đặc trưng của chèo, mà còn tạo nền tảng cho việc xây dựng hệ thống GraphRAG có khả năng truy vấn và suy luận hiệu quả trên dữ liệu văn hóa phi cấu trúc.
